\begin{statementbox}
\textbf{\textcolor{CStatement}{条件}}
\begin{description}[style=nextline, leftmargin=2.8em, labelsep=0.5em, itemsep=0.45em]
    \item[\textbf{\textcolor{CStatement}{条件一（标准椭圆）：}}]
    设椭圆 $C:\dfrac{x^2}{a^2}+\dfrac{y^2}{b^2}=1\ (a>b>0)$，
    中心为 $O(0,0)$，半焦距 $c$ 满足 $c^2=a^2-b^2$，离心率 $e=\dfrac ca$。
    \item[\textbf{\textcolor{CStatement}{条件二（点与斜率）：}}]
    取椭圆上非坐标轴点 $P(x_0,y_0)$，即 $x_0y_0\ne0$。
    记直线 $OP$ 的斜率为 $k_{OP}$，点 $P$ 处切线 $l$ 的斜率为 $k_l$。
\end{description}

\medskip
\textbf{\textcolor{CStatement}{核心结论}}
\begin{description}[style=nextline, leftmargin=2.8em, labelsep=0.5em, itemsep=0.45em]
    \item[\textbf{\textcolor{CStatement}{结论一（斜率乘积定值）：}}]
    \textbf{直观描述：} 从中心指向切点的直线与该点切线，斜率一正一负；
    两个斜率相乘，切点坐标消失，只留下椭圆的轴长比。
    \[
      \boxed{k_{OP}k_l=-\frac{b^2}{a^2}=e^2-1\in(-1,0).}
    \]
    \item[\textbf{\textcolor{CStatement}{结论二（考试速用）：}}]
    因 $k_{OP}=\dfrac{y_0}{x_0}\ne0$，故
    \[
      \boxed{k_l=-\frac{b^2}{a^2}\frac{x_0}{y_0}
             =\frac{e^2-1}{k_{OP}}.}
    \]
    反过来，若已知 $k_{OP}k_l$，则 $\boxed{e=\sqrt{1+k_{OP}k_l}}$。
    \item[\textbf{\textcolor{CStatement}{结论三（不依赖斜率的切线方程）：}}]
    对椭圆\emph{上任意一点} $P(x_0,y_0)$（包括四个顶点），均有
    \[
      \boxed{l:\ \frac{x_0x}{a^2}+\frac{y_0y}{b^2}=1.}
    \]
    此式适用范围比斜率乘积式更广。
\end{description}

\medskip
\textbf{\textcolor{CStatement}{与圆比较（特例而非退化）：}}
当 $a=b=R$，图形就是圆 $x^2+y^2=R^2$，此时 $e=0$。
在两个斜率都存在时 $k_{OP}k_l=-1$；本质是\textbf{圆的切线垂直于半径}。
在圆的坐标轴顶点，两条直线仍垂直，但斜率乘积没有定义。
\end{statementbox}
